\begin{prob}
    In discussion, we learned how to express the asymptotic time complexity of
    algorithms which depend on \emph{multiple} sizes. This problem will give more practice
    doing this.

    State the asymptotic time complexity of each of the operations below using
    multivariate $\Theta$ notation. Your answer should include every size variable used in
    the problem statement, such as number of points, dimensionality, etc. You
    do not need to show your work.

    \emph{Example}: compute the distance between a point in $\mathbb R^d$ and $n$ other
    points in $\mathbb R^d$.
    
    \emph{Example Solution}: $\Theta(nd)$.


    \begin{subprobset}
        \begin{subprob}
            Given a set of $n$ points in $\mathbb R^d$, compute the distance between
            each pair of points.

            \begin{soln}
                $\Theta(n^2 d)$.

                There are $\Theta(n^2)$ pairs, and computing the distance between a
                single pair of points takes $\Theta(d)$ time.
            \end{soln}
        \end{subprob}

        \begin{subprob}
            Given a set of $b$ blue points and a set of $r$ red points in $\mathbb R^d$,
            find the smallest distance between a blue point and a red point using the
            brute force approach.

            \begin{soln}
                $\Theta(brd)$.

                With a nested loop, we examine each pair where one point is red and the
                other is blue; there are $br$ such pairs. We calculate the distance for
                each pair in $\Theta(d)$ time.
            \end{soln}
        \end{subprob}


        \begin{subprob}
            Given a set of $n$ points in $\mathbb R^d$, 10\% of them red and
            90\% of them blue, compute the smallest distances between a blue point and a
            red point.

            \begin{soln}
                $\Theta(n^2 d)$.
            \end{soln}
        \end{subprob}

        \begin{subprob}
            Given a set $\mathcal D$ of $n$ data points in $\mathbb R^d$ and a
            new point $\vec x$ in $\mathbb R^d$, find the point in $\mathcal D$
            which is $k$th closest to $\vec x$ using the \python{kth_smallest}
            function below:

            \begin{minted}[autogobble]{python}
                def kth_smallest(numbers, k):
                    """Given an array of numbers, find index of kth smallest."""
                    numbers = np.copy(numbers)
                    for _ in range(k):
                        smallest_index = None
                        smallest_value = float('inf')
                        for i, x in enumerate(numbers):
                            if x < smallest_value:
                                smallest_index = i
                                smallest_value = x
                        numbers[smallest_index] = float('inf')
                    return smallest_index, smallest_value
            \end{minted}

            \begin{soln}
                $\Theta(n d + nk)$.

                For each point in the data set $\mathcal D$, we calculate its distance
                to $\vec x$ in time $\Theta(d)$. Doing this for all $n$ points takes
                $\Theta(nd)$ time. Finding the $k$th smallest takes $\Theta(nk)$ time.
                These two steps together take $\Theta(nd + nk)$ time.
            \end{soln}
        \end{subprob}

    \end{subprobset}

\end{prob}
